%% ===================================================================
%% Bootstrap recovery: liveness under validator loss + finality safety
%% Companion proof to LP-304 (Bootstrap Trust Anchor — weak-subjective
%% recovery decoupled from consensus finality).
%%
%% Establishes:
%%   (Thm 1) the prior model (bootstrap gated on >2/3 of CURRENT stake)
%%           is UNSATISFIABLE for mass recovery: there is a reachable
%%           system state from which no node can ever bootstrap.
%%   (Thm 2) the LP-304 BootstrapTrust policy is LIVE: a node converges
%%           whenever an honest configured-beacon supermajority is
%%           reachable, EVEN with > 1/3 of validators down.
%%   (Thm 3) the decoupling preserves FINALITY SAFETY: no block is
%%           finalized except by > 2/3 of current validator stake; the
%%           bootstrap anchor never finalizes (C1 preserved).
%% ===================================================================

\documentclass[11pt]{article}
\usepackage{amsmath,amssymb,amsthm}
\usepackage{booktabs}
\usepackage{enumitem}
\usepackage{hyperref}
\usepackage[margin=1in]{geometry}

\hypersetup{colorlinks=true,linkcolor=black,citecolor=blue,urlcolor=blue}

\newtheorem{theorem}{Theorem}
\newtheorem{lemma}[theorem]{Lemma}
\newtheorem{corollary}[theorem]{Corollary}
\newtheorem{definition}{Definition}
\newtheorem{proposition}[theorem]{Proposition}
\newtheorem{remark}{Remark}

\title{Bootstrap Recovery Liveness under Validator Loss,\\
       with Preserved Consensus-Finality Safety}
\author{Lux Industries --- Consensus Group \\ Companion proof to LP-304}
\date{2026-06-27}

\begin{document}
\maketitle

\begin{abstract}
We formalise the separation, introduced in LP-304, of the \emph{bootstrap trust} threshold
from the \emph{consensus finality} threshold. We model a primary network of weighted
validators a subset of which may be offline, including the very nodes attempting to recover.
We prove (i) the prior design, which named the bootstrap frontier by a quorum over
$>\!\tfrac23$ of \emph{current total stake}, has a reachable deadlock state in which no node
can bootstrap; (ii) the LP-304 policy, naming the frontier from an authenticated
configured-beacon supermajority with a response floor, is live whenever an honest
configured-beacon supermajority is reachable, even with more than $\tfrac13$ of validators
offline; and (iii) the two-threshold decoupling preserves BFT finality safety: a block is
accepted only under $>\!\tfrac23$ of \emph{current} validator stake, and the bootstrap anchor
never finalises. The third result is the weak-subjectivity soundness condition (``C1'')
under which the liveness gain costs no safety.
\end{abstract}

\section{Model}

\begin{definition}[Validator set and stake]
Let $V=\{v_1,\dots,v_n\}$ be the current primary-network validators with stake weights
$w:V\to\mathbb{Z}_{>0}$ and total $W=\sum_{v\in V}w(v)$. The two-thirds floor is
$\mathsf{F}(W)=\lfloor 2W/3\rfloor$; a stake mass $s$ \emph{clears} the floor iff $s>\mathsf{F}(W)$.
\end{definition}

\begin{definition}[Reachability]
At a point in time a subset $R\subseteq V$ is \emph{reachable} (connected and responding). A
node $u$ attempting to recover is itself offline-for-recovery: $u\notin R$. Write
$W_R=\sum_{v\in R}w(v)$ for the reachable stake.
\end{definition}

\begin{definition}[Configured beacon anchor]
Independent of $V$, each node holds a \emph{configured beacon set} $B$ resolved, in priority
order, from (1) operator configuration, (2) a trusted finalized checkpoint, or (3) genesis.
$B$ is fixed at the node and \textbf{not} derived from peer self-report. Let $B_R=B\cap R$ be
the reachable beacons and $\mathsf{honest}(B_R)\subseteq B_R$ the honest ones.
\end{definition}

\begin{definition}[Bootstrap policy, LP-304]\label{def:policy}
A $\mathsf{BootstrapPolicy}$ is $(\,B,\ \rho,\ m,\ \mathsf{cp}\,)$ with agreement ratio
$\rho\in\{\tfrac23,\tfrac34\}$, response floor $m\in\mathbb{Z}_{\ge 2}$, and optional
checkpoint $\mathsf{cp}$. On authenticated replies from responders $B_R$, it
\textbf{accepts} a frontier block $b$ iff
\[
   |B_R|\ \ge\ m \quad\text{and}\quad w\!\left(\{\,v\in B_R : b \preceq \mathsf{tip}(v)\,\}\right)\ >\ \rho\cdot W_{B_R},
\]
where $b\preceq \mathsf{tip}(v)$ means $b$ is the reported tip of $v$ or an ancestor of it
(ancestor-tolerant naming, LP-304 \S algorithm), and $W_{B_R}=\sum_{v\in B_R}w(v)$ is the
\emph{responding} beacon weight; or $\mathsf{cp}$ explicitly authorises $b$. The denominator
is the responding beacons, never $W$.
\end{definition}

\section{Unsatisfiability of the prior (current-stake) bootstrap threshold}

\begin{definition}[Prior threshold]
The prior design names a frontier only if the stake of validators agreeing on it clears
$\mathsf{F}(W)$ over the \emph{current total} $W$.
\end{definition}

\begin{theorem}[Mass-recovery deadlock]\label{thm:deadlock}
There is a reachable state in which the prior threshold can never be met, hence no offline
node bootstraps, even though the network is otherwise correct.
\end{theorem}

\begin{proof}
Take $n=5$, $w\equiv c$, so $W=5c$ and $\mathsf{F}(W)=\lfloor 10c/3\rfloor=3c$ (for $c$ a
multiple of $3$, and $3c<10c/3$ so the strict floor is $3c$). Let two validators be offline
for recovery, $R=\{v_1,v_3,v_4\}$, $W_R=3c$. The maximum stake that can agree on any frontier
is $W_R=3c$, and $3c \not> 3c$. The floor is never cleared, independent of agreement quality
or ancestor tolerance. The two offline validators contribute the missing $2c$; but they are
exactly the nodes whose recovery is blocked. The blocked set and the missing stake coincide,
so the state is a fixed point: recovery requires their stake, their stake requires their
recovery. $\qed$
\end{proof}

\begin{remark}
Theorem~\ref{thm:deadlock} is not an eclipse or a fault: $R$ is honest and reachable. It is a
threshold-model error --- the \emph{finality} denominator $W$ used for a \emph{liveness}
(bootstrap) decision. This is precisely the mainnet state observed for \texttt{luxd-2}
(LP-304, \S empirical): $W=2.5\times10^{18}$, $\mathsf{F}=1.\overline{6}\times10^{18}$,
$W_R=1.5\times10^{18}=60\%<66.\overline{6}\%$.
\end{remark}

\section{Liveness of the LP-304 policy}

\begin{theorem}[Recovery liveness under validator loss]\label{thm:liveness}
Let the policy be $(B,\rho,m,\cdot)$ with $\rho\le\tfrac34$. Suppose the reachable honest
beacons satisfy $|B_R|\ge m$ and $w(\mathsf{honest}(B_R))>\rho\cdot W_{B_R}$, and the honest
beacons agree (tip-or-ancestor) on a common block $b^\star$. Then the policy accepts a
frontier $b\succeq$ the honest-agreed height and the node syncs to it --- even if
$W_R\le \mathsf{F}(W)$, i.e.\ more than $\tfrac13$ of \emph{total} validators are offline.
\end{theorem}

\begin{proof}
The honest beacons all have $b^\star\preceq\mathsf{tip}(\cdot)$, so the agreeing weight on
$b^\star$ is at least $w(\mathsf{honest}(B_R))>\rho\,W_{B_R}$, clearing the policy's
\emph{responding-beacon} threshold; and $|B_R|\ge m$ clears the floor. By
Definition~\ref{def:policy} the policy accepts $b^\star$ (or a higher block on the same chain
with the same property; ancestor tolerance selects the highest such). Acceptance depends only
on $W_{B_R}$ and $|B_R|$, never on $W$; the offline validators' absence from $B_R$ does not
enter. Hence acceptance holds for any amount of offline \emph{total} stake, including
$W_R\le\mathsf{F}(W)$, as long as the honest reachable-beacon supermajority condition holds.
For the $n=5$, $R=\{v_1,v_3,v_4\}$ instance with $B=V$, $m=3$, $\rho=\tfrac23$: $|B_R|=3\ge 3$
and (all three honest) the agreeing weight $3c>\tfrac23\cdot 3c=2c$, so $b^\star=1082796$ is
accepted --- the case Theorem~\ref{thm:deadlock} deadlocks. $\qed$
\end{proof}

\begin{proposition}[Anti partition-capture]\label{prop:floor}
With $m\ge 2$ and the strict ratio $\rho$, an attacker who partitions the node to
$|B_R|<m$ reachable beacons cannot make the policy accept any frontier, regardless of their
unanimity.
\end{proposition}

\begin{proof}
Acceptance requires $|B_R|\ge m$. If $|B_R|<m$ the first conjunct of
Definition~\ref{def:policy} fails outright; unanimity ($\rho\,W_{B_R}$ trivially cleared)
is irrelevant. Only an explicit operator checkpoint $\mathsf{cp}$ can authorise a frontier
below the floor, by deliberate operator action. $\qed$
\end{proof}

\section{Finality safety is preserved (weak-subjectivity soundness, ``C1'')}

\begin{definition}[Finality]
A block is \emph{final} iff it carries a verified $\alpha$-of-$K$ quorum certificate whose
signing weight clears $\mathsf{F}(W)$ over the \emph{current} validator stake (LP-045/142).
$\mathsf{HasFinality}(s,W)\equiv s>\mathsf{F}(W)$.
\end{definition}

\begin{theorem}[Decoupling preserves safety]\label{thm:safety}
The LP-304 separation does not change the finality predicate: a block is accepted as final
iff $\mathsf{HasFinality}$ holds on current stake. The bootstrap policy never finalises a
block, and no frontier it names is treated as final.
\end{theorem}

\begin{proof}
$\mathsf{BootstrapTrust.AcceptsFrontier}$ returns a value of type $\mathsf{Frontier}$ (a sync
target), never a certificate, and the descent re-executes and re-verifies every block before
the node re-enters consensus (Invariant~3); after re-entry, acceptance is dispatched solely
through $\mathsf{ConsensusQuorum.HasFinality}$ on the current stake set, syntactically distinct
from $\mathsf{BootstrapTrust}$. The two interfaces share no implementation: the policy's
denominator is responding-beacon weight $W_{B_R}$, the finality denominator is current stake
$W$. Therefore for any block $b$ marked final there is a certificate with weight $>\mathsf{F}(W)$;
the bootstrap path produces no such certificate and cannot mark $b$ final. $\qed$
\end{proof}

\begin{theorem}[Weak-subjective soundness of the named frontier]\label{thm:c1}
Under the assumption that the configured beacon set contains less than a $\rho$-fraction of
Byzantine responding weight, the policy never names a forged (non-canonical) frontier.
\end{theorem}

\begin{proof}
A forged block $f$ is on $\mathsf{tip}(v)$'s chain only for Byzantine $v$ (an honest beacon
reports its real accepted chain, and the ancestor relation $f\preceq\mathsf{tip}(v)$ is
content-addressed: $f$'s id is the hash of its bytes, and the descent follows verified
$\mathsf{parent}$ links, so a substituted parent yields a different id and is unreachable ---
LP-304 \S rationale). Hence the agreeing weight on $f$ is at most the Byzantine responding
weight $<\rho\,W_{B_R}$, which does not clear the policy threshold; $f$ is never named (tests
D, E). The canonical block, backed by the honest $>\rho$ supermajority, is named instead.
An attacker cannot enlarge the Byzantine fraction by injecting beacons, since $B$ is resolved
from configuration/checkpoint/genesis and a non-configured peer has zero weight
(Invariant~1, test~D). $\qed$
\end{proof}

\section{Conclusion}

\begin{corollary}
Separating $\mathsf{BootstrapTrust}$ from $\mathsf{ConsensusQuorum}$ converts an
unsatisfiable recovery (Thm~\ref{thm:deadlock}) into a live one under realistic validator
loss (Thm~\ref{thm:liveness}, Prop~\ref{prop:floor}) at no cost to finality safety
(Thm~\ref{thm:safety}) and no new forgery surface under the standard weak-subjectivity
assumption (Thm~\ref{thm:c1}). Validators may freely leave and rejoin: a rejoining node
bootstraps from whichever configured beacons are reachable and agreeing, then finalises on
the current stake set. The liveness threshold and the safety threshold are now distinct
objects with distinct denominators, and cannot be re-conflated.
\end{corollary}

\end{document}
